\chapter{Komputasi yang Menghilang}
\label{ch:pervasif}

Pada musim dingin 2025/26 saya mengikuti dua kuliah pervasive computing dari Alois Ferscha. Yang
satu membahas sistem dan lingkungan: visi komputasi pervasif, sensor dan modalitas, jaringan sensor,
cara mengenali manusia dan benda, arsitektur pengenalan konteks, dan komputer yang dikenakan di
kepala. Yang lain membahas perancangan dari sisi manusia: persepsi visual, gerak mata, pendengaran,
perabaan, atensi, emosi, antarmuka yang bisa disentuh, dan gerakan motorik. Latihannya meminta kami
mengumpulkan data sensor sendiri, lalu membangun pengenal aktivitas dan sistem pelacak gerak.

\section{Visi Weiser}

\citet{weiser1991} membuka artikelnya dengan kalimat yang sering dikutip: teknologi yang paling
mendalam adalah teknologi yang menghilang, yang menenun dirinya ke dalam kehidupan sehari-hari sampai
tidak bisa dibedakan darinya. Menulis adalah contohnya. Kita tidak lagi memikirkan teknologi tulisan
ketika membaca papan nama jalan. Weiser membayangkan komputer yang sama: tersebar di dinding, di
meja, di pakaian, bekerja tanpa diminta, dan tidak menuntut perhatian. Hari ini ponsel, jam pintar,
dan rumah pintar mendekati visi itu, dengan semua manfaat dan masalahnya.

Komputasi pervasif membutuhkan tiga kemampuan: merasakan dunia lewat sensor, memahami konteksnya,
dan bertindak dengan cara yang tidak mengganggu. Bagian tengah, memahami konteks, adalah tempat
machine learning bekerja.

\section{Sensor dan modalitas}

Ponsel biasa membawa akselerometer, giroskop, magnetometer, mikrofon, kamera, GPS, dan sensor cahaya
serta jarak. Di luar ponsel, ada label RFID yang bisa dibaca tanpa kontak, pelacak mata, sensor
tekanan di lantai, dan sensor fisiologis seperti detak jantung dan konduktansi kulit. Setiap modalitas
punya kelebihan dan kekurangannya. Kamera kaya informasi tetapi mengganggu privasi dan boros energi.
Akselerometer murah dan hemat, tetapi sinyalnya sulit ditafsirkan. Kuliah menekankan satu hal: tidak
ada sensor yang langsung mengukur konteks. Tidak ada ``sensor rapat'' atau ``sensor sedang
memasak''. Konteks harus disimpulkan dari kombinasi sinyal mentah.

Jaringan sensor menambahkan masalah komunikasi dan energi. Simpul-simpul kecil bertenaga baterai harus
memutuskan kapan mengukur, kapan mengirim, dan kapan tidur. Memproses data di simpul itu sendiri,
alih-alih mengirim semuanya ke server, sering menghemat energi karena mengirim satu bit lewat radio
lebih mahal daripada banyak operasi hitung.

\section{Rantai pengenalan konteks}

Arsitektur pengenalan konteks yang dibahas kuliah mengikuti rantai yang juga dijelaskan dalam tutorial
\citet{bulling2014}. Data mentah dari sensor dipotong menjadi jendela-jendela waktu yang tumpang tindih,
misalnya dua detik dengan geseran satu detik. Pada frekuensi sampling 50 Hz, setiap jendela berisi
seratus sampel per sumbu. Dari setiap jendela dihitung fitur: rata-rata, simpangan baku, korelasi
antar sumbu, energi di pita frekuensi tertentu dari FFT, atau jumlah puncak. Fitur-fitur itu
dimasukkan ke pengklasifikasi, misalnya $k$-nearest neighbour, pohon keputusan, SVM, atau jaringan
saraf, yang menghasilkan label aktivitas untuk setiap jendela. Langkah terakhir menghaluskan urutan
label, karena seseorang tidak berganti dari berjalan ke duduk lalu ke berjalan lagi dalam tiga detik.

Rantai ini terlihat sederhana, tetapi kuliah dan latihan memperlihatkan jebakannya. Yang paling penting
adalah cara evaluasi. Kalau jendela-jendela dari orang yang sama masuk ke data latih dan data uji
sekaligus, model belajar mengenali gaya gerak orang itu, bukan aktivitasnya, dan skornya tampak
terlalu bagus. Evaluasi yang jujur meninggalkan satu orang di luar pelatihan, lalu mengujinya pada
orang itu, sebuah variasi dari pembagian data yang hati-hati di \cref{ch:data}. Dalam latihan kami,
yang dikenali adalah cara menggenggam dan gerakan memasang sekrup, direkam dengan akselerometer ponsel.
Setiap jendela diringkas menjadi sekitar empat puluh fitur, dari rata-rata, simpangan baku, kemiringan
distribusi, dan laju perpotongan nol di setiap sumbu, sampai korelasi antar sumbu, frekuensi dominan,
dan entropi spektral.

\section{Melacak gerak}

Sensor inersia memberi kecepatan sudut dari giroskop dan percepatan dari akselerometer. Mengintegrasikan
kecepatan sudut memberi orientasi, tetapi setiap bias kecil ikut terintegrasi dan orientasi perlahan
melenceng. Akselerometer, saat sensor diam atau bergerak pelan, mengukur arah gravitasi dan memberi
orientasi kemiringan yang tidak melenceng, tetapi berisik. Kedua sumber ini saling melengkapi.
Filter komplementer menggabungkannya dengan cara yang sederhana:
\begin{equation}
\theta_k=\alpha\big(\theta_{k-1}+\omega_k\Delta t\big)+(1-\alpha)\,\theta^{\mathrm{acc}}_k,
\end{equation}
dengan $\alpha$ dekat satu, misalnya $0{,}98$. Giroskop dipercaya untuk perubahan cepat, akselerometer
untuk kebenaran jangka panjang. Ini versi tangan dari Kalman filter di \cref{ch:prob}, dengan gain yang
dipilih sendiri alih-alih dihitung dari kovarians. Untuk melacak posisi, percepatan harus diintegrasikan
dua kali, dan galatnya tumbuh kuadratik terhadap waktu, sehingga pelacak posisi inersia murni hanya
berguna untuk waktu singkat kecuali ada sumber koreksi lain seperti kamera atau GPS.

\section{Manusia di dalam lingkaran}

Separuh kedua materi membahas sisi manusia, karena sistem pervasif berinteraksi langsung dengan
persepsi dan perhatian kita. Mata tidak bergerak mulus. Ia melompat dalam gerakan cepat yang disebut
sakade, lalu berhenti dalam fiksasi, dan hanya di daerah kecil di tengah penglihatan kita melihat
dengan tajam. Pelacak mata merekam pola ini, dan pola itu mengungkap apa yang menarik perhatian,
seberapa berat beban pikiran, atau apakah seseorang mengantuk. Memakai pandangan mata sebagai
input punya masalah klasik yang disebut masalah sentuhan Midas: mata kita melihat banyak hal tanpa
bermaksud memilihnya, sehingga sistem harus bisa membedakan melihat dari memilih.

Komputasi afektif mencoba mengenali emosi. Model emosi yang dipakai beragam, dari kategori emosi dasar
seperti senang, sedih, marah, dan takut, sampai ruang dua dimensi valensi dan gairah, di mana setiap
emosi adalah sebuah titik. Sinyalnya bisa berupa wajah, suara, postur, atau sinyal fisiologis. Kuliah
mengingatkan bahwa emosi yang diungkapkan tidak selalu sama dengan yang dirasakan, dan bahwa
pengenalan emosi menyentuh wilayah privasi yang pribadi.

Antarmuka yang bisa disentuh, atau \istilah{tangible user interface}, memberi bentuk fisik pada informasi
digital, sehingga kita bisa memegang, memindah, dan menyusun benda nyata untuk mengendalikan komputer.
Kuliah juga melihat ke depan, ke materi yang bisa diprogram dan mengubah bentuknya sendiri.

\section{Di mana machine learning masuk}

Hampir seluruh rantai pengenalan konteks kini bisa diganti jaringan saraf yang belajar fiturnya sendiri
dari sinyal mentah, seperti jaringan konvolusi satu dimensi atau LSTM pada data akselerometer
(\cref{ch:lstm}). Analisis adegan akustik mengenali tempat dan kejadian dari suara, dan pada prinsipnya
memakai jaringan visi pada spektrogram. Tantangan khas bidang ini tetap ada: data berlabel mahal karena
harus direkam dengan orang sungguhan, setiap orang bergerak berbeda, sensor bergeser posisinya, dan
perangkatnya punya baterai dan memori kecil. Model kecil yang bisa berjalan di mikrokontroler, sering
disebut TinyML, menjadi cabang tersendiri. Dan karena sensornya melekat di tubuh dan di rumah,
pertanyaan privasi dan persetujuan tidak bisa dipisahkan dari pertanyaan teknis (\cref{ch:manusia}).

\sumberbab{Bab ini mengikuti kuliah Pervasive Computing: Systems and Environments (visi, sensor,
jaringan sensor, identifikasi manusia dan benda, arsitektur pengenalan konteks, komputasi yang
dikenakan di kepala) dan Pervasive Computing: Design and Development (persepsi, pelacakan mata,
pendengaran, perabaan, atensi, komputasi afektif, antarmuka tangible, gerak motorik), beserta latihan
pengenalan aktivitas dan pelacakan gerak.}
